\documentclass[11pt,a4paper]{amsart}
\usepackage{amsmath,amssymb,amsthm,mathtools}
\usepackage[margin=1.2in]{geometry}
\usepackage{hyperref}

\theoremstyle{plain}
\newtheorem{theorem}{Theorem}[section]
\newtheorem{lemma}[theorem]{Lemma}
\newtheorem{proposition}[theorem]{Proposition}
\newtheorem{corollary}[theorem]{Corollary}
\theoremstyle{definition}
\newtheorem{definition}[theorem]{Definition}
\newtheorem{conjecture}[theorem]{Conjecture}
\theoremstyle{remark}
\newtheorem{remark}[theorem]{Remark}

\newcommand{\ZZ}{\mathbb{Z}}
\newcommand{\QQ}{\mathbb{Q}}
\newcommand{\CC}{\mathbb{C}}
\newcommand{\T}{T}
\DeclareMathOperator{\Res}{Res}

\title[A closed form for the Brown--Zudilin middle row]{A compact closed form for the
middle companion row\\ of the Brown--Zudilin $\zeta(5)$ construction}
\author{[Author]}
\date{}

\begin{document}
\begin{abstract}
Brown and Zudilin's cellular construction for $\zeta(5)$ produces three rational sequences
$Q_n,\widehat P_n,P_n$, all annihilated by one explicit order-three operator. Only $Q_n$ was
given by a closed binomial sum; $\widehat P_n$ and $P_n$ are defined by the recurrence. We prove
that the middle row has the two-term closed form
$\widehat P_n=\sum_{k,l}\T(n,k,l)\bigl(H^{(3)}_{n+k}-\Psi H^{(2)}_{n+k}\bigr)$, where $\T$ is the
cellular summand and $\Psi$ an explicit antisymmetric weight. The proof combines a Barnes--Mellin
descent with Zudilin's partial-fraction lemma, and uses no comparison of $\QQ$-coefficients of
zeta values --- a point of substance, since the required linear independence is not known. We
record what the same methods do \emph{not} give for the top row $P_n$, in the form of five
bounded exclusions. Two companion results are included for contrast: the corresponding
$\zeta(3)$ closed form, which is machine-verified; and the $\zeta(2)$ case, where the closed
form is known numerically but the method \emph{provably} fails, by an obstruction that is the
smallest instance of the one still blocking the top row.
\end{abstract}
\maketitle

\section{Introduction}

Brown and Zudilin \cite{BZ} construct, from cellular integrals on $\overline{\mathcal M}_{0,8}$,
a family of linear forms in $1,\zeta(2)\zeta(3),\zeta(5)$ whose coefficients are three rational
sequences $Q_n,\widehat P_n,P_n$. All three satisfy a single order-three recurrence with
polynomial coefficients, obtained in \cite{BZ} by creative telescoping on the integral. The
sequences are then \emph{defined} by that recurrence together with initial values; the agreement
of two order-three solutions at $n=0,1,2$ is what proves their decomposition.

Of the three, only $Q_n$ is given in \cite{BZ} by a closed binomial sum. Our object here is the
middle row.

\begin{theorem}\label{thm:main}
For every $n\ge0$,
\[
  \widehat P_n=\sum_{k,l=0}^{n}\T(n,k,l)\,\widehat w_3(n,k,l),
  \qquad \widehat w_3=H^{(3)}_{n+k}-\Psi\,H^{(2)}_{n+k},
\]
with $\T$, $\Psi$ as in Definition \ref{def:objects}.
\end{theorem}

Two features of the proof are worth flagging at the outset.

First, it is not a creative-telescoping certificate. The operator $L$ of \cite{BZ} is
\emph{not} a telescoper of $\T\cdot\widehat w_3$; the minimal telescoper of that summand has
order seven (\S\ref{sec:neg}). The proof instead runs through a Barnes--Mellin descent, where the
order-three structure lives.

Second, and less obviously, the proof must avoid comparing coefficients of zeta values. The
decomposition of \cite{BZ} identifies the $\zeta(2)$-coefficient of an integral with
$-4\widehat P_n$; a Barnes evaluation identifies it with something else; and concluding that the
two agree requires $\QQ$-linear independence of $1,\zeta(2),\zeta(3),\zeta(4),\zeta(5),
\zeta(2)\zeta(3)$, which is not known --- indeed the irrationality of $\zeta(5)$ is open.
Lemma \ref{lem:norm} is what removes that step.

\medskip
\noindent\textbf{Status.} Theorem \ref{thm:main} is proved. It is \emph{not} machine-verified;
\S\ref{sec:formal} states precisely what has been formalised and what has not. The
corresponding statement for the top row $P_n$ is open, and \S\ref{sec:neg} records what is
excluded.

\section{The objects}

\begin{definition}\label{def:objects}
Write $H^{(r)}_m=\sum_{j=1}^m j^{-r}$ and $H_m=H^{(1)}_m$. Put
\[
  \T(n,k,l)=\binom{n+k}{n}\binom nk^2\binom{n+l}{n}\binom nl^2\binom{n+k+l}{n},
  \qquad Q_n=\sum_{k,l=0}^{n}\T(n,k,l).
\]
For $r\ge1$ set
\[
  A_r(x)=H^{(r)}_{n+x}-H^{(r)}_x,\quad
  B_r(x)=H^{(r)}_{n-x}-H^{(r)}_x,\quad
  C_r=H^{(r)}_{n+k+l}-H^{(r)}_{k+l},
\]
and
\[
  \alpha=A_1(k)-A_1(l),\quad \beta=B_1(k)-B_1(l),\quad \Psi=\tfrac12\alpha+\beta,
  \qquad L_k=-A_1(k)-C_1-2B_1(k),
\]
with $L_l$ the mirror of $L_k$ under $k\leftrightarrow l$. Let $L$ denote the order-three
operator of \cite{BZ}, and let $\widehat P$ be the solution of $L\widehat P=0$ with
$\widehat P_0=0$, $\widehat P_1=\frac{101}4$, $\widehat P_2=\frac{344923}{96}$.
\end{definition}

Note $\alpha,\beta,\Psi$ are antisymmetric under $k\leftrightarrow l$, while $\T$ is symmetric.

\section{Two reductions}

The first is free, and is used twice more below.

\begin{lemma}\label{lem:sym}
For any weight $w$, $\sum_{k,l}\T\,w=\sum_{k,l}\T\,w^{\mathrm{sym}}$, where
$w^{\mathrm{sym}}(k,l)=\frac12\bigl(w(k,l)+w(l,k)\bigr)$. In particular every
$k\leftrightarrow l$-antisymmetric weight is annihilated by the double sum.
\end{lemma}

\begin{proof}
$\T(n,k,l)=\T(n,l,k)$, so relabelling gives $\sum\T\,w=\sum\T\,w^{\sigma}$ with
$w^\sigma(k,l)=w(l,k)$; average.
\end{proof}

Explicitly, writing $X_r=H^{(r)}_{n+k}$ and $Y_r=H^{(r)}_{n+l}$,
\begin{equation}\label{eq:w3sym}
  \widehat w_3^{\,\mathrm{sym}}=\tfrac12(X_3+Y_3)-\tfrac12\Psi\,(X_2-Y_2),
\end{equation}
since $\Psi$ is antisymmetric.

\begin{lemma}[the $\zeta(2)$ bridge]\label{lem:bridge}
For every $n\ge0$,
\[
  \sum_{k,l=0}^{n}\T\Bigl(H^{(3)}_{k+l}+\tfrac14(L_k+L_l)H^{(2)}_{k+l}\Bigr)
  \;=\;\sum_{k,l=0}^{n}\T\,\widehat w_3^{\,\mathrm{sym}}.
\]
\end{lemma}

\begin{proof}[Proof sketch]
Both sides arise as the $\zeta(2)$-coefficient of one Barnes-type double integral. Writing the
translated rational kernel as
$R(x,y)=P(x)P(y)P(x+y)/\bigl(Q(x)^2Q(y)^2\bigr)$ with $P(z)=\prod_{r=1}^n(z-r)$ and
$Q(z)=\prod_{r=0}^n(z+r)$, the local data at the double pole $y=-l$ is
$g_l(x)=c_lP(x)P(x-l)/Q(x)^2$, and $\sum_kB_{kl}=0$ for each $l$ by vanishing of the residue at
infinity. Since $L_k-L_l=-2\Psi$, symmetry reduces the difference of the two sides to
$-2\sum_l\sum_{j=l+1}^{n}g_l'(j)$; and on $l<j\le n$ the function $g_l$ has a \emph{double} zero,
because $P(x)$ vanishes at $x=j$ (as $1\le j\le n$) and $P(x-l)$ vanishes at $x=j$ (as
$1\le j-l\le n$), while $Q(j)\neq0$ for $j>0$. Hence every $g_l'(j)=0$.
\end{proof}

\begin{remark}
The vanishing ranges here are sharp: $g_l(j)=0$ exactly on $1\le j\le n+l$, $g_l'(j)=0$ exactly
on $l<j\le n$, and the simple-pole coefficient $q_l$ satisfies $q_l(j)=0$ exactly on
$1\le j\le n$ with a zero of order exactly one.
\end{remark}

\section{The descent, and its normalisation}

Specialising \cite[(I3)]{BZ} at $k=n+l$ writes the companion integral $I''_n$ as
\[
  I''_n=\sum_{l}(-1)^l\binom{n+l}{n}\binom nl^2\,J_3(n,n,n,n-l;n,n,n+l),
\]
and the rational part of the shifted $J_3$ is supplied by Zudilin's partial-fraction
lemma \cite[Lemma 4]{Zu02}: for a suitable one-variable $R$ with
$R(t)=\sum_kA_k(t+k)^{-2}+\sum_kB_k(t+k)^{-1}$, the induced quantity is
$2A\zeta(3)-B$, where $A=\sum_kA_k$ and
\[
  B=2\sum_kA_k\,H^{(3)}_{k-h}+\sum_kB_k\,H^{(2)}_{k-h}.
\]
For the specialisation at hand one computes, at $t=-n-1-k$,
\begin{equation}\label{eq:AB}
  A_{kl}=\binom{n+k}{n}\binom{n+k+l}{n}\binom nk^2,
  \qquad B_{kl}=A_{kl}\,L_k .
\end{equation}

The descent carries a normalisation, and the following lemma is what makes it harmless.

\begin{lemma}[normalisation]\label{lem:norm}
With the outer factor $(-1)^l\binom{n+l}{n}\binom nl^2$ of the specialisation,
\[
  \binom{n+l}{n}\binom nl^2\cdot A_{kl}=\T(n,k,l)
\]
identically in $n,k,l$. Consequently the $\zeta(3)$-coefficient produced by
\cite[Lemma 4]{Zu02} equals $2Q_n$, whereas that of $I''_n$ is $Q_n$; the descent therefore
computes $2I''_n$, and the factor is $2$ for \emph{every} $n$.
\end{lemma}

\begin{proof}
The two products consist of the same six binomial factors, reordered. The second assertion
follows since $Q_n\ge1$.
\end{proof}

\begin{remark}
This is the step that avoids any independence hypothesis. Both quantities being compared are
explicitly computed \emph{rational} coefficients inside a single derivation, not real numbers;
no statement about $\QQ$-linear independence of zeta values is used, here or anywhere else in
the proof.
\end{remark}

\begin{proposition}\label{prop:descent}
For every $n\ge0$,
$\displaystyle \widehat P_n=\sum_{k,l=0}^{n}\T\Bigl(H^{(3)}_{k+l}+\tfrac12L_kH^{(2)}_{k+l}\Bigr)$.
\end{proposition}

\begin{proof}
Combine the specialisation with \eqref{eq:AB}, \cite[Lemma 4]{Zu02} and Lemma \ref{lem:norm};
the outer sign cancels against the $(-1)^l$ of the $\zeta(3)$-coefficient.
\end{proof}

\begin{proof}[Proof of Theorem \ref{thm:main}]
By Lemma \ref{lem:sym} the summand of Proposition \ref{prop:descent} may be replaced by its
symmetrisation, which is $H^{(3)}_{k+l}+\frac14(L_k+L_l)H^{(2)}_{k+l}$. Lemma \ref{lem:bridge}
identifies its double sum with $\sum\T\,\widehat w_3^{\,\mathrm{sym}}$, and Lemma \ref{lem:sym}
again replaces that by $\sum\T\,\widehat w_3$.
\end{proof}

\section{What the same methods do not give}\label{sec:neg}

The top row remains open: we do not know a closed form for $P_n$. The following exclusions
delimit where it is not, and are all obtained by exact computation over stated ranges or by
proof, as indicated.

\begin{proposition}\label{prop:order7}
$L$ is not a telescoper of $\T\cdot\widehat w_3$. The minimal telescoper of that summand is
$A\cdot L$ with $A$ of order four, unique up to scalar, hence of order seven.
\end{proposition}

The corresponding certificate is far too large to be useful: its cofactors carry $1.5\cdot10^6$
monomials with coefficients of $280$ bits and more.

\begin{proposition}\label{prop:wstar}
There is a weight $w^\star$, explicitly a $29$-monomial element of the degree-two bare alphabet,
with $\sum\T\,w^\star=\widehat P_n$ and such that $L$ \emph{is} a telescoper of
$\T\cdot w^\star$.
\end{proposition}

Telescoper order is thus a property of the summand, not of the sum. This is what makes
Proposition \ref{prop:order7} compatible with the existence of an order-three certificate.

\begin{proposition}\label{prop:bridge}
No identity $\T\cdot(w^\star-\widehat w_3)=\Delta_kR_0+\Delta_lS_0$ with vanishing boundary terms
exists. Any operator $M$ with $M\cdot\T\cdot(w^\star-\widehat w_3)$ a discrete divergence must
annihilate $Q$; consequently the minimal such order is at least four.
\end{proposition}

\begin{proof}
For a monomial $M_i$ maximal in the shift closure the relevant shift matrices are diagonal, so
the $M_i$-component is a scalar equation; multiplying by $\T$ and summing over
$0\le k,l\le n$ telescopes the right-hand side to boundary terms, giving $d_iQ_n=0$ where
$d=w^\star-\widehat w_3$. Since $Q_n\ge1$, $d_i=0$; but $d$ has $29$ nonzero maximal components.
The second assertion follows from the same computation applied to $M$, together with the fact
that $Q$ has no annihilator of order one or two with polynomial coefficients of degree $\le13$.
\end{proof}

\begin{proposition}\label{prop:w5}
$L$ is not a telescoper of $\T\cdot w$ for any weight $w$ with $\sum\T\,w=P_n$ lying in the
nine-symbol bare weight-five span of degree at most three.
\end{proposition}

This is an exact computation: the admissible space has dimension $449$ of $1270$, and $P_n$
misses its image under the sum map by $292$ of $501$ equations. The obstruction is sharp --- no
single block family excludes, but two pairs do minimally --- and the same code returns the
positive answer of Proposition \ref{prop:wstar} at weight three when run as a control.

Finally, a negative about the literature route. Zudilin's higher-weight
statement \cite[Lemma 19]{Zu02} isolates $\zeta(5)$ for a one-variable very-well-poised $R$,
but does not specialise to the double kernel used here; and \cite{BZ} explicitly leaves the
decomposition of the two-variable universal kernels undone, supplying a descent for $I''$ only.
So the argument of \S4 has, at present, no weight-five analogue.

\section{The $\zeta(3)$ analogue, which is machine-verified}\label{sec:z3}

For contrast, and because the comparison is instructive, we record the corresponding statement
one weight down, where the formalisation \emph{is} complete.

Let $\mathcal A(n,k)=\binom nk^2\binom{n+k}k^2$, let $a_n=\sum_k\mathcal A(n,k)$ be the Ap\'ery
numbers, and let $b$ be the solution of Ap\'ery's recurrence
$(n+2)^3b_{n+2}=P(n+1)b_{n+1}-(n+1)^3b_n$, $P(n)=34n^3+51n^2+27n+5$, with $b_0=0$, $b_1=6$.

\begin{theorem}[machine-verified]\label{thm:z3}
For all $n\ge0$,
\[
  b_n=\sum_{k=0}^{n}\binom nk^2\binom{n+k}k^2\Bigl(2H^{(3)}_n-H^{(3)}_k\Bigr).
\]
\end{theorem}

\begin{theorem}[machine-verified]\label{thm:V}
For all $n\ge0$,
$\displaystyle\sum_{k=0}^{n}\binom nk^2\binom{n+k}k^2\bigl(H_{n+k}+H_{n-k}-2H_k\bigr)=0 .$
\end{theorem}

Theorem \ref{thm:V} has a purely rational proof: with
$\varphi(z)=\prod_{i=1}^n(z+i)\big/\prod_{j=0}^n(z-j)$, the simple poles of $\varphi$ have
residues $(-1)^k\binom nk\binom{n+k}k$ by Lagrange interpolation, and
$\Res_{z=m}\varphi^2=2\beta_m\sum_{k\ne m}\beta_k/(m-k)$ is antisymmetric in $(m,k)$, hence sums
to zero. No contour argument and no growth estimate is required.

Both are formalised in Lean 4 with no \texttt{sorry}, and depend on only the three
standard axioms (\texttt{propext}, choice, and quotient soundness). The weight
$2H^{(3)}_n-H^{(3)}_k$ appears to be new; we have not located it in the literature.

The scope of Theorem \ref{thm:z3} deserves the same care as \S\ref{sec:formal} demands of
Theorem \ref{thm:main}: in the formal development $b$ is pinned by the \emph{recurrence}
together with $b_0=0,b_1=6$, not by Ap\'ery's classical double sum with its nested alternating
inner sum. The equivalence of the two is proved on paper by six Zeilberger/Gosper certificates
and is not formalised. The situation is thus exactly parallel to the $\zeta(5)$ case, one step
further advanced.

\section{The $\zeta(2)$ case: where the method provably stops}\label{sec:z2}

The same template applied to the Ap\'ery-$\zeta(2)$ pair produces the closed form but
\emph{cannot} prove it, and the failure is sharp enough to locate exactly. We record it because
it is the smallest instance of the obstruction that also blocks the top row of
\S\ref{sec:neg}.

Let $S(n,k)=\binom nk^2\binom{n+k}k$, $A(n)=\sum_kS(n,k)$ (the sequence
$1,3,19,147,1251,\dots$), annihilated by $(n+1)^2u_{n+1}=(11n^2+11n+3)u_n+n^2u_{n-1}$, and put
\begin{equation}\label{eq:z2}
  B(n)=\sum_{k}S(n,k)\cdot\tfrac15\Bigl[H^{(2)}_n+H_k\bigl(2H_k-H_{n-k}-H_n\bigr)\Bigr].
\end{equation}
Then $L[B]=0$ holds for $n\le25$ by exact computation. We do not know a proof.

Everything in the argument of \S\ref{sec:z3} survives as far as the certificate. With
\[
  \rho(n,k)=k^2+k(1+6n)-4-15n-11n^2
\]
--- degree two, the smallest certificate we have met anywhere --- the Wilf--Zeilberger step
closes, and the structural boundary identity holds \emph{exactly}:
$\rho(n,n+1)=-2(n+1)(2n+1)=-(n+1)^2\binom{2n+2}{n+1}\big/\binom{2n}{n}$. The correction
decomposes as $A+B_1H_k+B_2(H_{n-k}+H_n)$ with $B_2=\rho(n,k+1)/\bigl(5(k+1)\bigr)$. Then it
stops.

\begin{proposition}\label{prop:z2}
The $H_{n-k}$ component of \eqref{eq:z2} is not Gosper-summable.
\end{proposition}

\begin{proof}
Write $T(k)=S(n,k)B_2(n,k)$. Then $T(k+1)/T(k)=\frac{a(k)}{b(k)}\cdot\frac{c(k+1)}{c(k)}$ with
$a(k)=(n-k)^2(n+k+1)$, $b(k)=(k+1)^2(k+2)$, $c(k)=\rho(n,k+1)$, and
$\gcd\bigl(a(k),b(k+j)\bigr)=1$ for all $j\ge0$, the roots being $n,-n-1$ against $-j-1,-j-2$
with $n$ generic. Gosper's equation is $a(k)x(k+1)-b(k-1)x(k)=c(k)$. Since
\[
  a(k)-b(k-1)=-n\bigl(k^2+(n+2)k-n(n+1)\bigr)
\]
has degree two with leading coefficient $-n\ne0$, the degree bound gives $\deg x=0$. Writing
$x=x_0$, the coefficient of $k^2$ forces $x_0=-1/n$, and the coefficient of $k$ then reads
$n+2$ against the required $3+6n$. Contradiction.
\end{proof}

Independently, the corresponding linear system has no solution for any $\deg x\le8$.

\begin{remark}\label{rem:z2}
The diagnosis is structural, not accidental. In the cases where the template does close --- the
$\zeta(3)$ pair of \S\ref{sec:z3}, and likewise Franel's numbers and $s_{10}$ --- the weight is
$k\leftrightarrow n-k$ symmetric, the certificate module has rank two, and each component
telescopes on its own. Here the weight is genuinely three-lettered in $H_k$, $H_{n-k}$, $H_n$;
the components do not telescope separately, and the weight-two part of any enlarged ansatz is
forced to vanish, since $\Delta(Sz)=0$ implies $z=0$. Closing this case therefore requires
either a different gauge for the weight --- the harmonic-monomial fit has a large kernel, and a
kernel shift changes the correction --- or a genuine $\Pi\Sigma$-ansatz admitting nested sums
rather than rational $\times$ harmonic monomials.

This is the same shape of obstruction as at weight five, where the top period
$\zeta(5)+2\zeta(2)\zeta(3)$ is impure in exactly the sense of carrying two independent letters.
The $\zeta(2)$ case is a miniature of it, small enough to be computed in full, and the two
available fixes are the same two.
\end{remark}

\section{Formal verification of the $\zeta(5)$ rows: what exists}\label{sec:formal}

Theorem \ref{thm:main} is a proof on paper. A formalisation in Lean 4 (\texttt{ZetaLucas}) is in
progress and is \emph{not} complete; we state its position precisely because the distinction is
easy to lose.

Verified, with no \texttt{sorry} and only the three standard axioms: a reflective
polynomial-identity checker; the summand shift calculus; the harmonic letter calculus in the
shifted base; the implication ``recurrence plus initial values implies the closed form''; and the
initial values of the relevant double sums, computed from their definitions.

Not verified: that either double sum satisfies $L$. Two hypotheses remain quarantined, and the
$\zeta(5)$ closed forms are consequently not machine-checked. Two further points are worth
recording. First, what the order-three certificate of Proposition \ref{prop:wstar} would give,
once complete, is $\widehat P_n=\sum\T\,w^\star$ --- \emph{not} Theorem \ref{thm:main}, since by
Proposition \ref{prop:bridge} the two representatives cannot be bridged at order zero. Second,
during this work a named hypothesis carried in the development was found to be false as stated;
every theorem depending on it was \texttt{sorry}-free with clean axioms, because
\texttt{\#print axioms} reports what a proof depends on and not whether its hypotheses are
inhabited. We mention it because the same trap is available to any development that reports
conditional results.

\begin{thebibliography}{9}
\bibitem{BZ} F.~Brown and W.~Zudilin, \emph{On cellular rational approximations to $\zeta(5)$},
arXiv:2210.03391.
\bibitem{Zu02} W.~Zudilin, \emph{Arithmetic of linear forms involving odd zeta values},
J.~Th\'eor.\ Nombres Bordeaux \textbf{16} (2004), 251--291.
\end{thebibliography}

\end{document}
